\begin{textsample}{NEG Dim 3 – pi – Score -3.00 – pi/pi\_em\_46.txt}  \label{ex:f3_neg_026}
A Área de Ciências da Natureza e a Integração entre os Componentes A integração entre os componentes curriculares da área de Ciências da Natureza e suas Tecnologias, no Ensino Médio, permite que os estudantes possam construir e utilizar conhecimentos específicos da área para argumentar, propor soluções e enfrentar desafios locais e/ou globais, relativos às condições de vida e ao ambiente.
Comprometendo -se, assim como as demais áreas, com a formação dos jovens para o enfrentamento dos desafios da contemporaneidade, na direção da educação integral e da formação cidadã, dando condições aos estudantes para aprofundar o exercício do pensamento crítico, realizar novas leituras do mundo, com base em modelos abstratos, e tomar decisões responsáveis, éticas e consistentes na identificação e solução de situações-problemas ( BRASIL, 2018 ).
A Área de Ciências da Natureza e suas Tecnologias traz um olhar articulado de seus componentes, para que as \textbf{competências} e habilidades se desenvolvam de forma conjunta, o que exige que o planejamento seja por área, indissociavelmente, e não por componente.
Nesse contexto, é apresentado o Currículo Referência da área de Ciências da Natureza e suas Tecnologias, instrumento norteador das estratégias de ensino, sistematização dos conteúdos e avaliação, proporcionando uma educação com uma abordagem crítica, caracterizando o conhecimento científico como cultural e relevante para a vida, proporcionando aos estudantes uma melhor compreensão e atuação no mundo contemporâneo.
Outro fator importante, para o ensino de Ciências da Natureza é que este possa favorecer a aprendizagem significativa do conhecimento histórico, o qual é acumulado na mente, para que se possa ter uma concepção de Ciência e suas relações com a Tecnologia e com a Sociedade.
Os campos do conhecimento científico – Astronomia, Biologia, Física, Geociências e Química – irão contemplar de forma multidisciplinar a área de Ciências da Natureza, buscando mudar a realidade das teorias vigentes, que se apresentam como conjuntos de proposições e metodologias altamente estruturadas e formalizadas, muito distantes do estudante em formação ( PCN, 1998, p. 27 ).
Na perspectiva de ensino por \textbf{competências} e habilidades, como é proposto pela BNCC, destacam -se indagações do tipo : o que são \textbf{competências}? e como aplicar estes novos conceitos de educação em sala de aula?
De acordo com Antunes ( 2001 ), \textbf{competências} são capacidades para resolver problemas ou para criar produtos que sejam considerados de valor em um meio social, são capacidades de compreender, de se adaptar e contextualizar.
No âmbito escolar, \textbf{competência} é a capacidade de mobilizar recursos, conhecimentos ou saberes vivenciados diante de situações complexas ( PERRENOUD, 2002 ).
De acordo com a BNCC, as \textbf{competências} são resultado da mobilização de conhecimentos, habilidades, valores e atitudes.
Desta forma, para se trabalhar por meio das \textbf{competências} e habilidades não significa que se deve deixar de ensinar conteúdos, mas sim deve -se, para além disso, utilizar informações, atribuindo -lhes um significado, fazendo uma contextualização com a vida e o espaço em que o estudante se insere.
Neste sentido, o professor deve procurar temas que se refletem no cotidiano dos alunos, para desta forma formular situações-problema.
Assim, os alunos serão desafiados e motivados a pesquisarem estas situações, a descobri -las e a apresentá -las ( BRASIL, 2018 ).
Além disso, segundo Lima et al., ( 2019 ), o Ensino das Ciências da Natureza tem o compromisso com o desenvolvimento do letramento científico, onde se entende que o estudante deve possuir a capacidade de realizar uma leitura do mundo, compreender e construir saberes e valores, e que possa atuar como um sujeito crítico capaz de identificar as múltiplas aplicações da ciência e da tecnologia no cotidiano.
Nesse contexto, o montante de informação disponível, permite aos professores encontrar uma maneira para que o aluno possa sair da passividade e desenvolver \textbf{competências}, se questionem e procurem conhecer sobre como produzir uma aprendizagem significativa e como construir o conhecimento ( PINTO et al., 2012 ).
Para isto, é necessário que o professor, durante a prática de ensino, insira diferentes metodologias, dando oportunidade aos alunos de externar seus conhecimentos prévios ( CARLOS \& BAHL, 2016 ).
Nesse âmbito, as metodologias ativas constituem -se formas de desenvolver o processo de aprender, colocando o aluno como sujeito ativo do seu processo de ensino e aprendizagem, utilizando experiências reais ou simuladas, com o intuito de estimular condições para solucionar desafios advindos das atividades essenciais da prática social, em diferentes contextos, ampliando possibilidades de exercitar a liberdade e a autonomia na tomada de decisões em diferentes momentos, preparando -se para o exercício profissional futuro ( BERBEL, 2011 ).

% matched lemmas: competência
\end{textsample}
